\documentclass[a4paper, 10pt]{article}
\usepackage{helvet}
\renewcommand{\familydefault}{\sfdefault}
\usepackage{pgf}
\usepackage{eurosym}
\usepackage{graphicx}
\usepackage{wasysym}
\usepackage{hyperref}
\usepackage{listings}
\usepackage{pxfonts}
\usepackage{verbatim}
\usepackage{color}
\usepackage{xcolor}
\usepackage{wrapfig}
\usepackage{enumitem}
\usepackage{booktabs}
\usepackage{gensymb}
\usepackage{tabularx}
\usepackage{currfile}

\hypersetup{
    bookmarks=true,         % show bookmarks bar?
    unicode=true,          % non-Latin characters in Acrobat’s bookmarks
    pdftoolbar=true,        % show Acrobat’s toolbar?
    pdfmenubar=true,        % show Acrobat’s menu?
    pdffitwindow=true,     % window fit to page when opened
    pdftitle={Assessments},    % title
    pdfauthor={Paul Vesey},     % author
    pdfsubject={I.C.T. Building Information Modelling},   % subject of the document
    pdfcreator={},   % creator of the document
    pdfproducer={xelatex}, % producer of the document
    pdfkeywords={'Graphics' }, % list of keywords
    pdfnewwindow=true,      % links in new PDF window
    colorlinks=true,       % false: boxed links; true: colored links
    linkcolor=violet,          % color of internal links (change box color with linkbordercolor)
    citecolor=magenta,        % color of links to bibliography
    filecolor=red,      % color of file links
    urlcolor=blue           % color of external links
}

\setlength\parindent{0pt}
\begin{document}

\lstset{language=HTML,
				basicstyle=\small,
				breaklines=true,
        numbers=left,
        numberstyle=\tiny,
        showstringspaces=false,
        aboveskip=-20pt,
        frame=leftline
        }
				
\begin{figure}
	\centering
	\includegraphics[width=0.5\linewidth]{./Assignments/img/LITlogo}
\end{figure}


\begin{tabularx}{\textwidth}{ |l|X| }
	\hline
	\textbf{Subject:} & COMP06051-ICT \& BIM\\
	\textbf{Course:} & BEng in Civil Engineering\\
	\textbf{Session:} & Autumn 2022\\
	\textbf{Lecturer:} & Paul Vesey \footnotesize{BEng, MIE, HDip}\\
	\textbf{Filename:} & \currfilebase\\
	\hline
\end{tabularx}



\vspace{0.25cm}	
	

\part*{Assignment 2 (20\%)- NBS Toolkit Project Definition; Group Project}


\begin{tabularx}{\textwidth}{ |X|X| }
	\hline
	\textbf{Issue Date:} & As stated on Moodle \\
	\hline 
	\textbf{Submission Date:}  & As stated on Moodle  \\
	\hline
\end{tabularx}

\section*{Introduction}

In this assignment you are required to utilise a Revit Model and associated specification to create an RIBA plan of work on the NBS toolkit.  The key deliverables for this project are:

\begin{enumerate}
	\item RIBA plan of work for Stages 4, 5, and 6 
	\item COBie Spreadsheets for each of the above stages
	\item Written report detailing your approach to the assignment and your evaluation of the toolkit
\end{enumerate}

\section*{Scenario}



In this project you are going to use the NBS toolkit to define a project in accordance with the RIBA Plan of Work. You are only required to detail Stages 4, 5, and 6. You do not have to complete Stages 0 to 3 or Stage 7.\\


This is a project for 2 individuals. All team members are expected to contribute equally to the project. As a team you are responsible for managing the workload and division of responsibility. All team members will receive the same mark, based on the academic merit of the submission.\\


The project is based on the Raheen Woods National School Project. The documents issued for tender have been uploaded to Moodle. You should examine these documents in detail to gain an understanding of the project.\\


The objective of the exercise is to create a comprehensive project management plan using the NBS toolkit. This will involve selecting appropriate tasks for each of the stages 4 through 6, developing deliverables for each stage, and ensuring the deliverables from each stage are incorporated into the next stage in a coherent manner.
\begin{itemize}
	\item Assume that Stage 4 starts on the 1\textsuperscript{st} of January 2019
	\item Assume that Stage 5 starts on the 1\textsuperscript{st} of July 2019
	\item Assume that Stage 6 handover will commence 30\textsuperscript{th} November 2019, and will be fully complete by the 1\textsuperscript{st} January 2020
\end{itemize}

\subsection*{Key Tasks for Each Stage}
\begin{enumerate}
	\item Provide summary Information
	\item Create Project Roles and assign fictitious individuals to those roles
	\item Add and delete tasks from the NBS Toolkit to reflect the needs of the project. Keep a note of the times deleted \& added, and explain the rationale for your decisions in a short report.
	\item Deliverables should be created in response to Project Management Tasks.
\end{enumerate}



\subsection*{Key Exemptions}

\begin{enumerate}
	\item You are not required to generate specifications for Deliverables as defined by the NBS Toolkit application. You are not required to use NBS-Generate
\end{enumerate}





\newpage
\section*{Recommended Approach}
Examine the documentation in detail. You have been provided with drawings and specifications.  The assignment involves creating a RIBA Stage 4 Plan, Stage 5 Plan and a Stage 6 plan.  You have, in effect, been provided with the output of Stage 4.
The NBS tool-kit is designed to capture data and assign responsibility. The key areas for each stage are:
\begin{enumerate}
	\item Details
	\item Rolls
	\item Tasks
	\item Deliverables
\end{enumerate}

The asset pack that accompanies this assignment are the deliverables of the Stage. You can use the items in the asset pack to create the deliverables of Stage 4. For instance, there is a detailed electrical specification for the building included in the asset pack. It naturally follows that a deliverable of Stage 4 is a ‘Detailed Electrical Specifications’. The associated Task will be the ‘Create Detailed Electrical Design and Specification. The next stage in reverse engineering Stage 4 will involve the creation of roles. So in this example, we can say that we need an electrical engineer; who will then be tasked with the design and specification of the electrical system.\\


This process is the essence of most project management; figure out what we need to do, what is involve in getting the task done and assigning the work to someone.\\



Stage 5 is about converting the detailed specification into real world assets. The NBS toolkit allows for the transfer (copy) of the deliverables between stages. This provides a starting point for the creation of the Stage 5 deliverables. They will have to be edited; for instance, the deliverable ‘Detailed Electrical Specification’ will have to be changed to ‘Installed Electrical System’ or similar. Again the same process can be followed to determine the tasks necessary to install the system, (install boards, switches, luminaries, pull cables etc.) and also to determine the rolls that will be required (Electrical Engineer, Supervisor, Electrician, etc.)  \\

Stage 6 will be completed in a similar fashion.\\



\newpage
\section*{Report}
Your report should be laid out in a professional manner, as would be done for a client or employer.  This involves, \emph{inter-alia}, a cover page, table of contents, and logical sectioning.  You report must include, but is not limited to:
\begin{itemize}
	\item Details of the methodology used to create each component
	\item Key findings and results.
	\item Details of any difficulties that you experienced extracting data from the applications
	\item Commentary on learning experience and the skills you developed during the production of this assignment.
\end{itemize}
  


\section*{Submission}
Your completed assignment comprising all computer files and written report are to be zipped into a single file and uploaded to Moodle on or before the date and time indicated.  Key components of the submission are:
\begin{itemize}
	\item NBS Toolkit RIBA plan of work Stage 4, 5, and 6.  I will require access to your plan of work on the system.  To do this you will need to invite me as a team member.  Use \href{mailto:paul.vesey@lit.ie}{paul.vesey@lit.ie} for the invitation.
	\item 3 off COBie Excel Spreadsheets
	\item Final report in pdf format.  Do not submit Word documents
\end{itemize}




\section*{Late Submission}
Failure to submit your assignment on or before the date and time indicated on Moodle will result in a penalty of 5\% per day or part thereof.

\subsection*{Marking Scheme}

\begin{table}[h!]
     \begin{center}
     \begin{tabular}{p{5cm}  p{5cm} }
     \toprule
      \textbf\large{Element} & \textbf\large{Proportion} \\ 
    \cmidrule(r){1-1}\cmidrule(lr){2-2}
      \textbf{RIBA Plan of Work} & \textbf{70\%}\\
      \textbf{Final Report} & \textbf{30\%}
      \\ \bottomrule
      \end{tabular}
      \label{tbl:markSchemeAsmt3}
      \end{center}
 \end{table}

\end{document}